\subsection{Homotopy-equivalent spaces}

DEFINITION 2. --- Let X and Y be topological spaces. A homotopy class \(\varphi \in [\mathrm{X};\mathrm{Y}]\) is invertible if there exists a homotopy class \(\psi \in [\mathrm{Y};\mathrm{X}]\) such that \(\psi \circ \varphi = [\mathrm{Id}_{\mathrm{X}}]\) and \(\varphi \circ \psi = [\mathrm{Id}_{\mathrm{Y}}]\). We say that a continuous map is a homotopy equivalence if its homotopy class is invertible.

Let \(\varphi\in[X;Y]\) be an invertible homotopy class. There exists a unique homotopy class \(\psi\in[Y;X]\) having the properties of Definition 2. Indeed, let \(\psi,\psi'\) be classes having these properties; we have

\[
\psi = \psi \circ [ \mathrm{Id} _ {\mathrm{Y}} ] = \psi \circ \varphi \circ \psi^ {\prime} = [ \mathrm{Id} _ {\mathrm{X}} ] \circ \psi^ {\prime} = \psi^ {\prime}.
\]

This unique class is called the inverse of the homotopy class \(\varphi\) and is denoted \(\varphi^{-1}\).

Let Z be a topological space. We have \(\chi\circ\varphi\circ\varphi^{-1}=\chi\) for all \(\chi\in[Y;Z]\) and \(\theta\circ\varphi^{-1}\circ\varphi=\theta\) for all \(\theta\in[X;Z].\) It follows that the maps \(\chi\mapsto\chi\circ\varphi\) from {[}Y;Z{]} to {[}X;Z{]} and \(\theta\mapsto\theta\circ\varphi^{-1}\) from {[}X;Z{]} to {[}Y;Z{]} are inverse bijections of each other. Likewise, the map \(\chi\mapsto\varphi\circ\chi\) from {[}Z;X{]} to {[}Z;Y{]} is bijective and its inverse is the map \(\theta\mapsto\varphi^{-1}\circ\theta\) from {[}Z;Y{]} to {[}Z;X{]}.

Let X, Y, Z be topological spaces, and let \(\varphi \in [X;Y]\) and \(\psi \in [Y;Z]\) be invertible homotopy classes. Then the class \(\psi \circ \varphi\) is invertible, with inverse \(\varphi^{-1} \circ \psi^{-1}\). Indeed, we have

\[
\begin{array}{r l} & (\psi \circ \varphi) \circ (\varphi^ {- 1} \circ \psi^ {- 1}) = \psi \circ (\varphi \circ \varphi^ {- 1}) \circ \psi^ {- 1} \\ & \qquad = \psi \circ [ \mathrm{Id} _ {\mathrm{Y}} ] \circ \psi^ {- 1} \\ & \qquad = \psi \circ \psi^ {- 1} = [ \mathrm{Id} _ {\mathrm{Z}} ] \end{array}
\]

and, similarly, \((\varphi^{-1} \circ \psi^{-1}) \circ (\psi \circ \varphi) = [\mathrm{Id}_{\mathrm{X}}]\).

Let X and Y be topological spaces and let \(f: X \to Y\) be a homotopy equivalence. Every continuous map \(g: Y \to X\) whose class is the inverse of that of f is said to be a homotopy inverse (or inverse) of f up to homotopy. Such a map g is a homotopy equivalence.

A homeomorphism \(f\) is a homotopy equivalence, and \([f]^{-1} = [f^{-1}]\). The composite of two homotopy equivalences is a homotopy equivalence. The relation “X and Y are topological spaces and there exists a homotopy equivalence from X to Y” is an equivalence relation.

DEFINITION 3. --- Topological spaces X and Y are said to be homotopy equivalent if there exists a homotopy equivalence from X to Y.

Examples. --- 1) The empty topological space is homotopy equivalent only to itself.

\begin{enumerate}
\def\labelenumi{\arabic{enumi})}
\setcounter{enumi}{1}
\item
  Let X be a nonempty topological space. For X to be homotopy equivalent to a space reduced to a point, it is necessary and sufficient that there exist a constant map \(p: X \rightarrow X\) which is homotopic to the identity map of X. (It is then so for every constant map \(q \colon X \to X\) since \([p] = [p] \circ [q] = [\mathrm{Id}_{\mathrm{X}}] \circ [q] = [q]\).) Indeed, let \(P\) be a space reduced to a point, \(f \colon P \to X\) a map, and \(g \colon X \to P\) the unique map from X to P. We have \(g \circ f = \mathrm{Id}_{\mathrm{P}}\); for f to be a homotopy equivalence, it is necessary and sufficient that \(f \circ g\) be homotopic to \(\mathrm{Id}_{\mathrm{X}}\). Now, \(f \circ g\) is constant, with image \(f(P)\).
\item
  A pointed topological space (X,x) is said to be contractible, or the topological space X is said to be contractible at x, if there exists a pointed homotopy at x connecting \(Id_{X}\) to the constant map from X into X with image \(\{x\}\). Such a space is homotopy equivalent to a point. There exist, however, spaces homotopy equivalent to a point which are contractible at none of their points, and spaces contractible at one point, but not at every point (III, p. 321, exerc. 1).
\item
  Let E be the n-dimensional Euclidean space (or, more generally, a topological vector space over R) and let X be a subset of E. The set X is said to be star-shaped at a point x of X if, for every \(y \in X\) and every \(t \in I\), the point \(tx + (1 - t)y\) belongs to X. A convex subset (I, p. 122) of E is star-shaped at each of its points.
\end{enumerate}

A topological subspace X of E that is star-shaped at one of its points x is contractible at that point. Indeed, the map \(\sigma\colon X \times I \to X\) defined by \(\sigma(y,t) = tx + (1 - t)y\) is a pointed homotopy at x connecting \(Id_{X}\) to the constant map with image \(\{x\}\).

In particular, every interval of R, every convex subset of a Euclidean space or, more generally, of a topological vector space over R, is contractible at each of its points.

\begin{enumerate}
\def\labelenumi{\arabic{enumi})}
\setcounter{enumi}{4}
\tightlist
\item
  We will demonstrate later (TA, V) that Euclidean spheres \(S_n\), \(n\geqslant1\), are not homotopy equivalent to a point. The unit sphere of a separable infinite-dimensional Hilbert space is contractible at each of its points (EVT, V, p. 71, exerc. 13).
\end{enumerate}

